\scanpage{534}
\mspage{445}
% archive sheet number: 533

Or on a le

\uline{Lemme} 4.15.7.4. Soient $C$ une catégorie, $x_0$ un objet de $C$, $p$ un projecteur de $x_0$, $\tau$ un endomorphisme de $x_0$ satisfaisant les relations

\begingroup\makeatletter\tagsleft@true\makeatother
\[\tau p=p\tau=\tau,\tag*{$\unsure{(a)|(*)}$}\]
\endgroup

et tel que l'endomorphisme de $x_0^\natural\overset{\mathrm{déf}}{=}\operatorname{Im}p$ (défini p. ex. comme objet de $C^\wedge$, $x_0^\natural=\operatorname{Ker}(\mathrm{id}_{x_0},p)$) soit l'identité, Alors on a $\tau=p$.

\uline{Démonstration} \inserted{de 4.15.7.4.} On est ramené aussitôt, par le \unsure{procédé} \unsure{standart|standard}, au cas où $C=\mathrm{Ens}$, où la vérification est immédiate : les endom. $\tau$ de l'ens $\deleted{\unsure{x|F}}E$ \inserted{$(=x)$} qui satisfont $(a)$ \deleted{sont \unsure{encore|aussi} de la forme} \inserted{correspondent \unsure{biunivoqmt|biunivoquemt}} aux \deleted{\unsure{de|le}} endom. de $E^\natural=pE$, via

\[\left\{\begin{aligned}u&\longmapsto\sigma u\rho\\ \tau&\mathrel{\unsure{\mapstochar\relbar\mapstochar|\longmapsto}}\tau\mid E^\natural\end{aligned}\right.\]

(où \diagram{\begin{tikzcd}[ampersand replacement=\&] E^\natural \arrow[r,hook,shift left,"\sigma"] \& E \arrow[l,shift left,"\rho"] \end{tikzcd}} \unsure{sont les flèches can.|est la flèche can.})

donc $\tau\mid E^\natural=\mathrm{id}\ \Longrightarrow\ u=\sigma\rho=p$.

Choisissons un $j_0\in J$, et considérons les hom.

\[\tau_{j,j_0}^\natural:x_0^\natural\xrightarrow{\sim}x_0^\natural\]

induits par les $\tau_{j,j_0}$ avec $j\geq j_0$. Ici encore, comme l'ens des hom. \deleted{\unsure{poss. en question|pour lesquels}} est de
